\section*{Reading Map}

\begin{itemize}[leftmargin=*,itemsep=1pt]
  \item \textbf{Part I: Access and Demand (Sections 1 to 8):} Why product excellence alone does not create commercial viability; the structural distinctions between promotion, distribution, and transactions; demand capture versus demand creation; the 4 forms of availability; and Harish Narayanan's who/what/how architecture.
  \item \textbf{Part II: Buying Journeys and Channel Architecture (Sections 9 to 17):} Channel families; transaction-value economics; the speed-versus-scale matrix; owned versus rented access; founder-hour efficiency; evolving channel-market fit; search intent layers; and the battle for distribution access.
  \item \textbf{Part III: Founder-Led and Owned Distribution (Sections 18 to 25):} Founder outreach as market apprenticeship; the ethics and mechanics of warm introductions; factual personalization; deep content as a product; content distribution loops; search reality in the age of platform volatility; authentic communities; and event-led growth.
  \item \textbf{Part IV: Paid Acquisition, Economics, and Attribution (Sections 26 to 37):} Paid media as experimentation; matching CAC to realized unit contribution; worked funnel conversion arithmetic; cash CAC versus economic CAC; the marginal CAC escalation curve; contribution ROAS; non-linear attribution; creative-destination coherence; retargeting limits; delivery capacity constraints; and the 3 marketing ledgers.
  \item \textbf{Part V: Partnerships, Channels, and Intermediaries (Sections 38 to 45):} Activation plans versus ceremonial MoUs; partner profit-and-effort models; partner readiness; taxonomy of intermediaries; account-based distribution; event classifications; creator reach versus peer trust; and platform risk stress tests.
  \item \textbf{Part VI: Cross-Industry Channel Systems (Sections 46 to 55):} Distribution across SaaS, CPG, professional services, enterprise software, higher education, local services, 2-sided marketplaces, luxury goods, industrial commodities, and new category creation.
  \item \textbf{Part VII: The Bayes Institute Go-to-Market Design (Sections 56 to 68):} Attention-constrained distribution (15 to 25 weekly founder hours); mapping 3 customer routes; the 1st canonical evidence asset; a 30-day D2C pilot; pilot economics; college partnership discovery; enterprise diagnostic wedge; prospect consent governance; weekly time budgets; paid-test criteria; cumulative content assets; and the initial growth loop.
  \item \textbf{Part VIII: The Founder's Distribution Operating System (Sections 69 to 76):} The 1-page channel thesis; the 10-dimension channel scorecard; the 13-line acquisition ledger; experimental stopping rules; brand investments versus demand tests; monthly portfolio reviews; the 90-day protocol; and 12 distribution traps.
  \item \textbf{Part IX: Advanced Channel Economics and Decision Mathematics (Sections 77 to 93):} Marginal cost analysis; capital allocation across channels; worked break-even tests; queueing constraints in sales funnels; double-counting attribution; audience saturation ceilings; AI answer engines; creative audience filtering; partner economics; brand versus demand time constants; adverse selection; seasonality; geographic channel shifts; the red-team exercise; and the advanced channel memo.
  \item \textbf{Appendices A to I and Final Synthesis:} Operational field templates, economic worksheets, audit checklists, comparative go-to-market systems, sensitivity tables for capacity cliffs, and the bridge to Chapter 21.
\end{itemize}

\vspace{4pt}

\section*{Part I: Distribution Is the Business of Access}

\subsection*{1. The Best Product Still Needs a Route to the Customer}

A brilliant educational curriculum can fail because qualified learners never discover its existence. A mediocre course can triumph because its brand is ubiquitous and its checkout flow is frictionless. A useful enterprise software tool can disappear beneath an incumbent that controls operating-system integrations. A manufacturer can build a superior domestic appliance and fail because it is absent from retail shelves at the point of purchase. None of these outcomes implies that product quality is irrelevant. Rather, they demonstrate that every commercial organization must produce 2 complementary goods: the intrinsic value customers consume, and a reliable, economically sustainable route through which that value is accessed.

The founder's 1st strategic distribution question is not ``Which digital platform is currently popular?'' It is: \textbf{Where does the customer's actual decision take place?} An engineering manager decides on developer tooling during an architecture review. An ambitious professional decides on specialized training after struggling with an unfamiliar production failure. A parent chooses educational support on the recommendation of a trusted peer. A procurement committee only considers suppliers pre-approved on an enterprise vendor master list. Distribution is the deliberate design of access under these binding operational conditions.

\subsection*{2. Promotion, Distribution, and Transaction Are Distinct Functions}

Modern digital interfaces collapse promotion, distribution, and transactions into a single click, misleading founders into equating an impression with distribution, and access with commercial revenue. In reality, these 3 functions operate on different mechanics:
\begin{itemize}[leftmargin=*,itemsep=1pt]
  \item \textbf{Promotion communicates:} It informs, persuades, or signals relevance. A technical video on YouTube creates awareness, but without an immediate, accessible pathway to trial, it produces zero commercial access.
  \item \textbf{Distribution provides practical access:} It places the offering at the buyer's choice point. A retail store that stocks an appliance in a locked backroom has formal inventory distribution, but lacks practical retail availability.
  \item \textbf{Transactions complete exchange:} They execute agreements, collect payments, and confirm contractual rights. An enthusiastic introductory meeting with a university dean represents goodwill, but it does not constitute a funded procurement transaction.
\end{itemize}

A complete distribution architecture guides an eligible prospect from relevant encounter to credible evaluation to feasible acquisition. It resolves practical friction: application workflows, syllabus downloads, payment rails, corporate billing, regional delivery, technical compatibility, and cohort scheduling. A founder who optimizes advertising impressions while ignoring these friction points is measuring activity instead of distribution.

\subsection*{3. Demand Capture and Demand Creation Solve Different Problems}

\textbf{Demand capture} intercepts buyers who are actively seeking a solution. Examples include an engineer searching Google for ``production RAG evaluation course with code review,'' a chief financial officer asking colleagues for a forecasting consultant, or a homeowner browsing energy-efficient ceiling fans. The commercial objective is to secure inclusion in an existing consideration set. Search-engine optimization, comparison directories, procurement participation, peer referrals, and high-intent keyword auctions are the primary tools of demand capture.

\textbf{Demand creation} shifts how buyers perceive and prioritize an unaddressed problem before they search for suppliers. A rigorous white paper demonstrating why naive language-model agents break in production can make evaluation training suddenly necessary. A diagnostic rubric illustrating that inventory stockouts stem from flawed forecasting can make an operational advisory service urgent. Original technical research, public teaching, and open-source benchmarks create vocabulary for latent pain. Demand creation operates over longer horizons, carries diffuse attribution, and requires patient capital. Founders must not mistake public intellectual appreciation for immediate commercial demand.

\subsection*{4. Demand May Be Latent but Commercially Inaccessible}

A market may harbor immense latent dissatisfaction yet remain commercially impenetrable. Thousands of mid-career software engineers may worry about artificial intelligence automation, yet lack the discretionary evening hours required for a demanding technical practicum. An enterprise may acknowledge severe inefficiencies in data pipelines, yet lack a designated budget owner with procurement authority to contract an external specialist. A university may admire modern machine-learning electives, yet its curriculum committees may have locked course schedules for the next 18 months. Distribution cannot conjure budget, executive authority, or customer time out of thin air.

The foundational unit of commercial distribution is the \textbf{buying situation}: a specific human being, an observable operational trigger, a recognized task, accessible alternative solutions, authorized financial resources, and a feasible purchasing route. This is substantially narrower than a demographic audience. One million technical practitioners on social media do not constitute 1 million reachable buyers. Sustainable distribution identifies precise moments when an acute operational problem converts into an authorized, funded commitment.

\subsection*{5. Distribution Has 4 Forms of Availability}

Drawing upon the empirical market foundations established by Byron Sharp and the Ehrenberg-Bass Institute, commercial access requires 4 distinct categories of availability:
\begin{enumerate}[leftmargin=*,itemsep=1pt]
  \item \textbf{Mental availability:} The probability that a prospective buyer retrieves, recognizes, and considers the organization when an acute operational need arises.
  \item \textbf{Physical or practical availability:} The practical ease with which the product or service can be acquired, inspected, scheduled, or integrated at the relevant moment.
  \item \textbf{Institutional availability:} The compliance of the seller with the buyer's internal policies, vendor onboarding protocols, data-governance standards, and procurement rails.
  \item \textbf{Economic availability:} The condition where total acquisition costs: including subscription fees, onboarding friction, time allocation, and organizational switching costs: remain acceptable to the buyer.
\end{enumerate}

An offering may possess high mental availability through widespread media coverage, yet suffer zero practical availability due to filled capacity or prohibitive scheduling. A specialized training vendor may possess technical excellence, yet fail institutional availability by lacking vendor registration. Channel design must systematically resolve the exact availability barrier that halts the transaction. Frequently, an operational modification: such as adding evening cohort slots or enabling corporate invoice billing: unlocks far more revenue than doubling advertising spend.

\subsection*{6. The Unit of Distribution Is a Customer Decision, Not a Platform Impression}

Consider an educational webinar attended by 2,000 participants. If 1,950 attendees are full-time university students seeking free slide decks, while the commercial product is an advanced ₹40,000 practicum designed for senior backend engineers, the webinar produced immense promotional reach but almost zero distribution access. Conversely, an intimate technical working group of 20 engineering managers who oversee production machine-learning deployments constitutes a potent distribution environment.

The governing denominator of distribution is always \textbf{qualified buyers encountered}, never impressions, views, or mailing-list counts. Founders must log who engaged, the context that triggered their interest, the observable operational need, the subsequent action taken, and the resulting economic contribution. A channel is commercially viable only when it reliably generates qualified prospects with demonstrated willingness to commit, at a customer acquisition cost that unit economics can sustainably bear.

\subsection*{7. Buying Processes Are Non-Linear}

Prospective buyers rarely navigate tidy conceptual funnels. An engineer may 1st discover Bayes Institute through a colleague's code review, search an unfamiliar technical term 3 months later, read a diagnostic essay, defer action, attend an open seminar, and finally enter the institute's URL directly to purchase. An enterprise buyer may circulate a white paper among 4 department heads, pause for budget renewals, and request a diagnostic 6 months after the initial touchpoint.

Last-click digital attribution routinely credits the final search ad or direct visit, disregarding the foundational assets that created conviction. Similarly, self-reported discovery surveys often capture the most recent or memorable interaction rather than the causal origin. Organizations must build \textbf{decision maps}: structured analyses of the triggers that initiate exploration, the authorities consulted, the alternatives evaluated, the proof required, the budget authorizers involved, and the administrative mechanisms that complete the sale. Funnels remain helpful for diagnosing stage-gate conversion losses; they must never be mistaken for the true psychological and institutional history of a sale.

\subsection*{8. Channel Choice Follows the Who, Why, and Promise}

In the Blume Ventures \textit{GTM Field Notes} (B2C Edition), Harish Narayanan establishes the primacy of the ``who/what/how'' sequence: founders must define the exact target buyer (who), their underlying operational need (why), and the precise product promise (what) before evaluating marketing channels (how). Selecting channels based on industry fads: such as launching paid search ads or short-form video campaigns simply because a consumer venture succeeded there: is a primary driver of startup capital destruction.

This principle is vital for complex, high-consequence offerings. When an institute offers advanced technical training to engineers responsible for critical production infrastructure, the optimal channels are specialized developer communities, engineering leadership roundtables, and technical failure post-mortems. Broad consumer ad networks may deliver inexpensive impressions, but the traffic lacks prerequisites, context, and purchasing urgency. Product-channel fit is not an abstract preference for a technology platform; it is a structural match between buyer context, problem urgency, communication depth, and unit economics.
